\chapter{Scripting Commands}
First of all, the Tcl interpreter created by CFS offers all the
standard commands known from e.g. {\em tclsh} or other Tcl
interpreters. For a detailed list have a look at
\url{http://www.tcl.tk/man/tcl8.4/TclCmd/contents.htm}.

Additionally, there is one new Tcl command available for interacting with
CFS, which has the general form
\begin{verbatim}
cfs <arguments>
\end{verbatim}
where \verb#arguments# is list of arguments, specifying the detailed
action and additional parameters. The second argument specifies the
recipient within CFS, e.g. the grid object or a pde, which offer
different commands. These can be found in the following sections or by
using the general command
\begin{verbatim}
cfs <object-name> getCommands
\end{verbatim}
which will return a Tcl list with all available functions of this
object including their parameters.

\section{Grid Object}
The object {\em grid} is accessible in every event procedure. It
provides access to the FE mesh, i.e. regions, elements and nodes.

\subsection{General Grid Information}

\tcl-command{getNumNodes}
Returns the total number of nodes in the grid.
\begin{tcl-params}
  \returnlist
  \param{numNodes}{unsigned int}{Number of nodes}
\end{tcl-params}

\tcl-command{getNumElems}
Returns the total number of elements in the grid.
\begin{tcl-params}
  \returnlist
  \param{numElems}{unsigned int}{Number of elements}
\end{tcl-params}

\tcl-command{getNumSurfElems}
Returns the number of surface elements in the grid.
\begin{tcl-params}
  \returnlist
  \param{numElems}{unsigned int}{Number of surface elements}
\end{tcl-params}

\tcl-command{getNumVolElems}
Returns the number of volume elements in the grid.
\begin{tcl-params}
  \returnlist
  \param{numElems}{unsigned int}{Number of volume elements}
\end{tcl-params}

\tcl-command{getNumElemsOfRegion}
Returns the number of elements of specified region.
\begin{tcl-params}
  \paramlist
  \param{regionName}{string}{Name of region}
  \returnlist
  \param{numElems}{unsigned int}{Number of elements}
\end{tcl-params}

\tcl-command{getNumNodesOfRegion}
Returns the number of nodes of specified region.
\begin{tcl-params}
  \paramlist
  \param{regionName}{string}{Name of region}
  \returnlist
  \param{numNodes}{unsigned int}{Number of nodes}
\end{tcl-params}

\tcl-command{getRegionNames}
Returns a list with the names of all regions (surface and volume)
contained in the grid.
\begin{tcl-params}
  \returnlist
  \param{regionNames}{list string}{List with region names}
\end{tcl-params}

\subsection{Node Access Functions}

\tcl-command{getNodeNames}
Returns the names of all named nodes.
\begin{tcl-params}
  \returnlist
  \param{names}{list string}{Names of all named nodes}
\end{tcl-params}

\tcl-command{getNodesByName}
Returns the node numbers of a nodelist with given name.
\begin{tcl-params}
  \paramlist
  \param{nodeName}{string}{Name of named nodes}
  \returnlist
  \param{nodeList}{list unsigned int}{Node numbers of named nodes}
\end{tcl-params}

\tcl-command{getNodesByRegion}
Returns the node numbers of all nodes contained in specified region.
\begin{tcl-params}
  \paramlist
  \param{nodeName}{string}{Name of region}
  \returnlist
  \param{nodeList}{list unsigned int}{Node numbers of given region}
\end{tcl-params}

\tcl-command{getNodeCoordinate}
Returns the coordinate of a node.
\begin{tcl-params}
  \paramlist
  \param{nodeNr}{unsigned int}{name of NodeList} 
  \returnlist
  \param{coordinate}{list double}{Coordinates \{x y z\} of node}
\end{tcl-params}


\subsection{Element Access Functions}

\tcl-command{getElemNames}
Returns the names of all named elements.
\begin{tcl-params}
  \returnlist
  \param{names}{list string}{Names of all named elements}
\end{tcl-params}



 
\section{PDE Objects}
The {\em pde} object refers to an abstract object responsible for
defining the physical field(s) to be calculated. Therefore it is
problem dependent and must be on of
\begin{itemize}
\item acoustic
\item electrostatic
\item magnetic
\item mechanic
\item thermal
\item direct
\end{itemize}
The methods of this object can be used mainly in the procedures
\verb#CFS_ReadBCs#, \verb#CFS_SetBCs# and \verb#CFS_PostProcess#.

\subsection{Result access}

\tcl-command{getValue}
Returns the nodal unknown(s) of a specified node.
\begin{tcl-params}
  \paramlist
  \param{nodeNr}{unsigned int}{Node number of node of interest}
  \returnlist
  \param{result}{list double}{Result of specified node}
\end{tcl-params}

\subsection{Boundary Conditions / Loads}

\tcl-command{setIdbc}
Allows one so set the value of (previously defined) inhomogeneous
Dirichlet boundary nodes, identiefied by their name.
\begin{tcl-params}
  \paramlist
  \param{name}{string}{Name of nodes}
  \param{value}{double}{Inhom. Dirichlet boundary value}
  \param{dof}{string}{Degree of freedom to be set (ux,uy,uz...)}
\end{tcl-params}

